% GENERATED FILE. Edit canonical lessons, then run npm run generate:books.
% canonical-source-hash: fnv1a64:9aa594a83b1edf77
% canonical-lessons: HI-C467-kapti, HI-C467-dhurt, HI-C467-bhola, HI-C467-visvasghat, HI-C467-makkar, HI-R467-review-460-467

\chapter{Honest and False}
\label{ch:hi-honest-and-false}
\hlchaptermodality{467}

\begin{chapteropening}
\textbf{By the end of this chapter:} \emph{I can say whether someone can be trusted.}

Complete the last lesson of chapter 467: I can say whether someone can be trusted.
\end{chapteropening}

\section[\texorpdfstring{\hi{कपटी} (kapṭī)}{कपटी (kapṭī)}]{\hi{कपटी} (kapṭī) — deceitful; two-faced}

\label{lesson:HI-C467-kapti}



\begin{quote}
Before the new one: say the Hindi for a gentleman, then the Hindi for harsh.
\end{quote}

\subsection*{You'll want to know: \hi{कपटी}}
\textbf{\hi{कपटी}} — \emph{kapṭī} — \textquotedblleft{}deceitful; two-faced\textquotedblright{}.

\begin{quote}
\hi{वह} \hi{आदमी} \hi{मीठा} \hi{बोलता} \hi{है}, \hi{लेकिन} \hi{बहुत} \hi{कपटी} \hi{है।}
\end{quote}

\begin{quote}
\emph{vah ādmī mīṭhā boltā hai, lekin bahut kapṭī hai.}
\end{quote}

\begin{quote}
That man talks sweetly, but he is very two-faced.
\end{quote}

\begin{cousinweb}
\textbf{\hi{कपटी}} (\emph{kapṭī}) is a tatsama, Sanskrit \emph{kapaṭin}, the adjective of \emph{kapaṭa}, \textquotedblleft{}fraud, deceit\textquotedblright{} (Monier-Williams).
\end{cousinweb}

\begin{grammarlens}[title={What you've built}]
One more word for describing what you lived and felt.
\end{grammarlens}

\subsection*{Your turn}
\begin{itemize}
\raggedright
  \item \emph{Say it:} \emph{kapṭī}
  \item \emph{Say it:} \emph{kapṭī}, once more
  \item \emph{Say it:} \emph{vah ādmī mīṭhā boltā hai, lekin bahut kapṭī hai.} — That man talks sweetly, but he is very two-faced.
  \item \emph{Recall:} say the Hindi for a gentleman, then the Hindi for harsh, then say \emph{kapṭī} again
\end{itemize}

\begin{tcolorbox}[breakable,colback=teal!4,colframe=teal!35!black,title={Before you move on}]
What does \hi{कपटी} mean? (\textquotedblleft{}Deceitful; two-faced\textquotedblright{}.) Say it once more.
\end{tcolorbox}
\section[\texorpdfstring{\hi{धूर्त} (dhūrt)}{धूर्त (dhūrt)}]{\hi{धूर्त} (dhūrt) — crafty; a cheat}

\label{lesson:HI-C467-dhurt}



\begin{quote}
Before the new one: say the Hindi for harsh, then the Hindi for deceitful.
\end{quote}

\subsection*{You'll want to know: \hi{धूर्त}}
\textbf{\hi{धूर्त}} — \emph{dhūrt} — \textquotedblleft{}crafty; a cheat\textquotedblright{}.

\begin{quote}
\hi{धूर्त} \hi{आदमी} \hi{ने} \hi{बूढ़ी} \hi{औरत} \hi{से} \hi{उसके} \hi{पैसे} \hi{ले} \hi{लिए।}
\end{quote}

\begin{quote}
\emph{dhūrt ādmī ne būṛhī aurat se uske paise le liye.}
\end{quote}

\begin{quote}
A crafty man took the old woman's money from her.
\end{quote}

\begin{cousinweb}
\textbf{\hi{धूर्त}} (\emph{dhūrt}) is a tatsama, borrowed from Sanskrit \emph{dhūrta}, \textquotedblleft{}cunning, crafty\textquotedblright{}, and as a noun \textquotedblleft{}a rogue, a cheat\textquotedblright{} (Monier-Williams).
\end{cousinweb}

\begin{grammarlens}[title={What you've built}]
One more word for describing what you lived and felt.
\end{grammarlens}

\subsection*{Your turn}
\begin{itemize}
\raggedright
  \item \emph{Say it:} \emph{dhūrt}
  \item \emph{Say it:} \emph{dhūrt}, once more
  \item \emph{Say it:} \emph{dhūrt ādmī ne būṛhī aurat se uske paise le liye.} — A crafty man took the old woman's money from her.
  \item \emph{Recall:} say the Hindi for harsh, then the Hindi for deceitful, then say \emph{dhūrt} again
\end{itemize}

\begin{tcolorbox}[breakable,colback=teal!4,colframe=teal!35!black,title={Before you move on}]
What does \hi{धूर्त} mean? (\textquotedblleft{}Crafty; a cheat\textquotedblright{}.) Say it once more.
\end{tcolorbox}
\section[\texorpdfstring{\hi{भोला} (bholā)}{भोला (bholā)}]{\hi{भोला} (bholā) — guileless; trusting}

\label{lesson:HI-C467-bhola}



\begin{quote}
Before the new one: say the Hindi for deceitful, then the Hindi for crafty.
\end{quote}

\subsection*{You'll want to know: \hi{भोला}}
\textbf{\hi{भोला}} — \emph{bholā} — \textquotedblleft{}guileless; trusting\textquotedblright{}.

\begin{quote}
\hi{मेरा} \hi{भाई} \hi{बहुत} \hi{भोला} \hi{है}, \hi{हर} \hi{आदमी} \hi{की} \hi{बात} \hi{मान} \hi{लेता} \hi{है।}
\end{quote}

\begin{quote}
\emph{merā bhāī bahut bholā hai, har ādmī kī bāt mān letā hai.}
\end{quote}

\begin{quote}
My brother is very trusting; he believes whatever anyone tells him.
\end{quote}

\begin{cousinweb}
\textbf{\hi{भोला}} (\emph{bholā}) is a tadbhava: CDIAL lists Hindi \emph{bholā}, \textquotedblleft{}foolish, guileless\textquotedblright{}, under \emph{*bhulla}, a form it reconstructs (CDIAL 9539). In Hindi it is mostly kind: simple, trusting, without tricks.
\end{cousinweb}

\begin{grammarlens}[title={What you've built}]
One more word for describing what you lived and felt.
\end{grammarlens}

\subsection*{Your turn}
\begin{itemize}
\raggedright
  \item \emph{Say it:} \emph{bholā}
  \item \emph{Say it:} \emph{bholā}, once more
  \item \emph{Say it:} \emph{merā bhāī bahut bholā hai, har ādmī kī bāt mān letā hai.} — My brother is very trusting; he believes whatever anyone tells him.
  \item \emph{Recall:} say the Hindi for deceitful, then the Hindi for crafty, then say \emph{bholā} again
\end{itemize}

\begin{tcolorbox}[breakable,colback=teal!4,colframe=teal!35!black,title={Before you move on}]
What does \hi{भोला} mean? (\textquotedblleft{}Guileless; trusting\textquotedblright{}.) Say it once more.
\end{tcolorbox}
\section[\texorpdfstring{\hi{विश्वासघात} (viśvāsghāt)}{विश्वासघात (viśvāsghāt)}]{\hi{विश्वासघात} (viśvāsghāt) — a betrayal of trust}

\label{lesson:HI-C467-visvasghat}



\begin{quote}
Before the new one: say the Hindi for crafty, then the Hindi for guileless.
\end{quote}

\subsection*{You'll want to know: \hi{विश्वासघात}}
\textbf{\hi{विश्वासघात}} — \emph{viśvāsghāt} — \textquotedblleft{}a betrayal of trust\textquotedblright{}.

\begin{quote}
\hi{दोस्त} \hi{का} \hi{विश्वासघात} \hi{अरुण} \hi{कभी} \hi{नहीं} \hi{भूलेगा।}
\end{quote}

\begin{quote}
\emph{dost kā viśvāsghāt Aruṇ kabhī nahīṃ bhūlegā.}
\end{quote}

\begin{quote}
Arun will never forget his friend's betrayal.
\end{quote}

Masculine.

\begin{cousinweb}
\textbf{\hi{विश्वासघात}} (\emph{viśvāsghāt}) is a tatsama, borrowed from Sanskrit \emph{viśvāsa-ghāta}, \textquotedblleft{}destruction of confidence, violation of trust\textquotedblright{} (Monier-Williams): \emph{viśvās}, \textquotedblleft{}trust\textquotedblright{}, and \emph{ghāt}, \textquotedblleft{}a blow\textquotedblright{}.
\end{cousinweb}

\begin{grammarlens}[title={What you've built}]
One more word for describing what you lived and felt.
\end{grammarlens}

\subsection*{Your turn}
\begin{itemize}
\raggedright
  \item \emph{Say it:} \emph{viśvāsghāt}
  \item \emph{Say it:} \emph{viśvāsghāt}, once more
  \item \emph{Say it:} \emph{dost kā viśvāsghāt Aruṇ kabhī nahīṃ bhūlegā.} — Arun will never forget his friend's betrayal.
  \item \emph{Recall:} say the Hindi for crafty, then the Hindi for guileless, then say \emph{viśvāsghāt} again
\end{itemize}

\begin{tcolorbox}[breakable,colback=teal!4,colframe=teal!35!black,title={Before you move on}]
What does \hi{विश्वासघात} mean? (\textquotedblleft{}A betrayal of trust\textquotedblright{}.) Say it once more.
\end{tcolorbox}
\section[\texorpdfstring{\hi{मक्कार} (makkār)}{मक्कार (makkār)}]{\hi{मक्कार} (makkār) — crooked; a cheat}

\label{lesson:HI-C467-makkar}



\begin{quote}
Before the new one: say the Hindi for guileless, then the Hindi for a betrayal of trust.
\end{quote}

\subsection*{You'll want to know: \hi{मक्कार}}
\textbf{\hi{मक्कार}} — \emph{makkār} — \textquotedblleft{}crooked; a cheat\textquotedblright{}.

\begin{quote}
\hi{वह} \hi{दुकानदार} \hi{मक्कार} \hi{है}, \hi{हमेशा} \hi{कम} \hi{तोलता} \hi{है।}
\end{quote}

\begin{quote}
\emph{vah dukāndār makkār hai, hameśā kam toltā hai.}
\end{quote}

\begin{quote}
That shopkeeper is a crook; he always gives short weight.
\end{quote}

\begin{cousinweb}
\textbf{\hi{मक्कार}} (\emph{makkār}) is usually traced to Arabic \emph{makkār}, \textquotedblleft{}crafty, wily\textquotedblright{}, which reached Hindi through Persian.
\end{cousinweb}

\begin{grammarlens}[title={What you've built}]
One more word for describing what you lived and felt.
\end{grammarlens}

\subsection*{Your turn}
\begin{itemize}
\raggedright
  \item \emph{Say it:} \emph{makkār}
  \item \emph{Say it:} \emph{makkār}, once more
  \item \emph{Say it:} \emph{vah dukāndār makkār hai, hameśā kam toltā hai.} — That shopkeeper is a crook; he always gives short weight.
  \item \emph{Recall:} say the Hindi for guileless, then the Hindi for a betrayal of trust, then say \emph{makkār} again
\end{itemize}

\begin{tcolorbox}[breakable,colback=teal!4,colframe=teal!35!black,title={Before you move on}]
What does \hi{मक्कार} mean? (\textquotedblleft{}Crooked; a cheat\textquotedblright{}.) Say it once more.
\end{tcolorbox}
\section[(dialogue)]{Looking back: from Missing Someone to Honest and False}

\label{lesson:HI-R467-review-460-467}



\begin{quote}
Say the Hindi for the pain of separation and to cling tightly.
\end{quote}

\subsection*{Your turn}
Say these aloud:

\begin{itemize}
\raggedright
  \item the pain of separation, to cling tightly, to grow restless, desolate, a longing
  \item ashamed, to shrink back shyly, to hesitate, to blush, brazen
  \item self-esteem, to put on airs, a boast, to strut, fussiness
  \item peace of mind, contentment, reassurance, to eat one's fill, consolation
  \item to get bored, interest, eager, curiosity, interesting
  \item a joke, to make someone laugh, to make someone cry, a sigh, wailing
  \item heartless, cruel, good-hearted, a gentleman, harsh
  \item deceitful, crafty, guileless, a betrayal of trust, crooked
\end{itemize}

\begin{tcolorbox}[breakable,colback=teal!4,colframe=teal!35!black,title={Before you move on}]
The pain of separation? (\textbf{\hi{विरह}}, \emph{virah}.) To cling tightly? (\textbf{\hi{चिमटना}}, \emph{cimaṭnā}.) To grow restless? (\textbf{\hi{अकुलाना}}, \emph{akulānā}.) Desolate? (\textbf{\hi{सूना}}, \emph{sūnā}.) A longing? (\textbf{\hi{लालसा}}, \emph{lālsā}.) Ashamed? (\textbf{\hi{शर्मिंदा}}, \emph{śarmindā}.) To shrink back shyly? (\textbf{\hi{सकुचाना}}, \emph{sakucānā}.) To hesitate? (\textbf{\hi{हिचकिचाना}}, \emph{hickicānā}.) To blush? (\textbf{\hi{लजाना}}, \emph{lajānā}.) Brazen? (\textbf{\hi{ढीठ}}, \emph{ḍhīṭh}.) Self-esteem? (\textbf{\hi{अभिमान}}, \emph{abhimān}.) To put on airs? (\textbf{\hi{इतराना}}, \emph{itrānā}.) A boast? (\textbf{\hi{शेख़ी}}, \emph{śekhī}.) To strut? (\textbf{\hi{अकड़ना}}, \emph{akaṛnā}.) Fussiness? (\textbf{\hi{नख़रा}}, \emph{nakhrā}.) Peace of mind? (\textbf{\hi{सुकून}}, \emph{sukūn}.) Contentment? (\textbf{\hi{संतोष}}, \emph{santoṣ}.) Reassurance? (\textbf{\hi{तसल्ली}}, \emph{tasallī}.) To eat one's fill? (\textbf{\hi{अघाना}}, \emph{aghānā}.) Consolation? (\textbf{\hi{सांत्वना}}, \emph{sāntvanā}.) To get bored? (\textbf{\hi{ऊबना}}, \emph{ūbnā}.) Interest? (\textbf{\hi{दिलचस्पी}}, \emph{dilcaspī}.) Eager? (\textbf{\hi{उत्सुक}}, \emph{utsuk}.) Curiosity? (\textbf{\hi{जिज्ञासा}}, \emph{jijñāsā}.) Interesting? (\textbf{\hi{रोचक}}, \emph{rocak}.) A joke? (\textbf{\hi{मज़ाक़}}, \emph{mazāq}.) To make someone laugh? (\textbf{\hi{हँसाना}}, \emph{hãsānā}.) To make someone cry? (\textbf{\hi{रुलाना}}, \emph{rulānā}.) A sigh? (\textbf{\hi{आह}}, \emph{āh}.) Wailing? (\textbf{\hi{विलाप}}, \emph{vilāp}.) Heartless? (\textbf{\hi{निष्ठुर}}, \emph{niṣṭhur}.) Cruel? (\textbf{\hi{क्रूर}}, \emph{krūr}.) Good-hearted? (\textbf{\hi{नेक}}, \emph{nek}.) A gentleman? (\textbf{\hi{सज्जन}}, \emph{sajjan}.) Harsh? (\textbf{\hi{कठोर}}, \emph{kaṭhor}.) Deceitful? (\textbf{\hi{कपटी}}, \emph{kapṭī}.) Crafty? (\textbf{\hi{धूर्त}}, \emph{dhūrt}.) Guileless? (\textbf{\hi{भोला}}, \emph{bholā}.) A betrayal of trust? (\textbf{\hi{विश्वासघात}}, \emph{viśvāsghāt}.) Crooked? (\textbf{\hi{मक्कार}}, \emph{makkār}.)
\end{tcolorbox}
